\section{Results}\label{sec:results}

\subsection{CRAM's benchmarks}\label{sec:cram}

We use CRAM's data, splits, test sets and operation count\cite{hoover2026classical}. On a grid, the check may use the
relations that cell coordinates express: same row, same column, same block for a block size $b$, the diagonals, and
conjunctions of two of these. The fit picks among them. This vocabulary is what lets one fitted rule serve every size;
CRAM instead learns a free table for each of the 3,240 pairs of cells of a Sudoku board.

\textbf{Sudoku.} On each of CRAM's \nm{sudoku.nfits} training splits of \nm{sudoku.trainboards} solved boards, the fit
keeps exactly the \nm{sudoku.kept} pair types that make up ``same row'', ``same column'' and ``same box'', and nothing
else. It takes \nm{sudoku.fit.lo} to \nm{sudoku.fit.hi} milliseconds on one CPU core. The rule solves all
\nm{sudoku.ntest} test puzzles and never stops at a wrong board (Table~\ref{tab:cram}). It uses \nm{sudoku.ops}
operations a puzzle under CRAM's convention, \nm{sudoku.ops.ratio} times fewer than CRAM, and \nm{sudoku.opscmp} with
every comparison counted. The cold stage alone solves \nm{sudoku.cold.pct}\% of the puzzles; a depth-first search
finishes the rest after a median of \nm{sudoku.moves.median} choices. The same rule, unchanged, solves
\nm{size4.solved} of \nm{size4.n} puzzles of $4\times4$, all by the cold stage, \nm{size16.solved} of \nm{size16.n} of
$16\times16$, and \nm{size25.final} of \nm{size25.n} of $25\times25$, with no wrong answer (Fig.~\ref{fig:puzzles}).
CRAM fits one model per board size.

\begin{table}[tbp]\centering\small
\caption{\figlabel{Table 1}{CRAM's benchmarks.} Operations a puzzle follow CRAM's convention, under which comparisons
are free; in brackets, every comparison counted as well. The rows for CRAM, its hand-set energy and TRMP, the most
accurate looped model, are from CRAM's Table~2 and Appendix~F\cite{hoover2026classical}. Our fit times are on one CPU
core, CRAM's on one H200 GPU. Ours on Sudoku-Extreme: each of three fits.}\label{tab:cram}
\begin{tabular}{llrrr}\toprule
Benchmark & Model & Solved & Operations a puzzle & Fit \\\midrule
Sudoku-Extreme (\nm{sudoku.ntest}) & TRMP, looped & \nm{trmp.pct}\% & \nm{trmp.ops} & \nm{trmp.fit} h \\
 & CRAM & \nm{cram.sudoku.pct}\% & \nm{cram.sudoku.ops} & \nm{cram.sudoku.fit} ms \\
 & CRAM, one weight set by hand & 100\% & \nm{cram.handset.ops} & none \\
 & \textbf{ours} & \textbf{100\%} & \nm{sudoku.ops} (\nm{sudoku.opscmp}) & \nm{sudoku.fit.lo}--\nm{sudoku.fit.hi} ms \\\midrule
Maze-Unique (\nm{unique.n}) & CRAM & 100\% & \nm{cram.unique.ops} & \nm{cram.unique.fit} ms \\
 & \textbf{ours} & \textbf{100\%} & \nm{unique.ops} (\nm{unique.opscmp}) & \nm{unique.fit} ms \\\midrule
Maze-Hard (\nm{hard.n}) & CRAM, hand-built distance layer & \nm{cram.hard.pct}\% & \nm{cram.hard.ops} & \nm{cram.hard.fit} \\
 & \textbf{ours}, no hand-built layer & \textbf{100\%} & \nm{hard.ops} (\nm{hard.opscmp}) & \nm{hard.fit} ms \\\bottomrule
\end{tabular}
\end{table}

\textbf{Mazes.} On Maze-Unique the fitted rule says that an open cell on the path has exactly two path neighbours, and
the start and the goal one; the cold stage alone solves all \nm{unique.n} test mazes, and it is dead-end filling,
which we did not write in. The same rule is right on all \nm{edge.n} tree mazes of \nm{edge.sides} cells a side whose
dead ends reach the board's edge, on which CRAM's released model is right on none (Section~\ref{sec:cramweights}). On
Maze-Hard, which has loops and labels one particular shortest path, we replace CRAM's hand-built distance layer by the
cold run's own clock: the loop at which a cell is first reached from the goal, a register that every looped cold block
has. The fitted rule is to step to a neighbour reached one loop earlier and, among several, to prefer down, then left,
right and up. That order is the benchmark's own tie-break, learned from the labels, not set by us. The rule is right on
all \nm{hard.n} test mazes, including the one that CRAM's model misses on every seed. It is also right on all
\nm{rooms.n} mazes of rooms and corridors that we built to probe CRAM's model (Section~\ref{sec:cramweights}), and on
\nm{loops.exact} of \nm{loops.n} mazes with loops of 31, 61 and 121 cells a side.

\textbf{N-Queens and colouring.} Fitted on the 92 solutions of the $8\times8$ board, the same procedure keeps the costs
for two queens on a row, a column or a diagonal. It completes \nm{queens.n} of \nm{queens.n} boards at $80\times80$ and
at $100\times100$, whether nine queens are missing, seven eighths of them or all of them. Fitted on planted
3-colourable graphs of 10 vertices, it keeps only ``the same colour across an edge'' and colours
\nm{colour.solved.3}, \nm{colour.solved.45} and \nm{colour.solved.6} of 100 graphs of 100 vertices at average degrees
3, 4.5 and 6. These instances are ours, not CRAM's test sets.

\subsection{CLRS-30}\label{sec:clrs}

We fit each task on \nm{clrs.train.inputs} inputs of the official training split, with no hints: 20 at each of the
sizes 4, 7, 11, 13 and 16, or, for the string matchers, at the library's one training length of 20 items. We score with
the library's scorer on the official test at 64 nodes. For larger sizes we use our own generator, which reproduces every
official test split we use bit for bit and agrees with all \nm{gen.labels} labels stored in the library's files for
\nm{gen.tasks} of the tasks; held-out inputs are drawn at seeds never used before.

\textbf{All thirty tasks are exact} on the official test and on every held-out set (Table~\ref{tab:clrs}): the official
validation split, fresh inputs at 64 nodes, and fresh graphs of 256 and \nm{held.max.nodes} nodes, dense as the library
draws them and sparse with about four neighbours a node; arrays and strings go to \nm{held.max.arrays} and
\nm{held.max.items} items. That is \nm{held.sets} held-out sets with no miss (Fig.~\ref{fig:clrs}).

Fifteen tasks need nothing but comparisons, counts and the completion of the answer: the four sorts, minimum, binary
search, quickselect, activity selection, task scheduling, breadth-first search, the minimum spanning trees of Kruskal
and Prim, longest common subsequence and the naive and Knuth--Morris--Pratt string matchers. Their completions are the
depth along the answer's pointer, the bottleneck closure of the weights and the closure of agreeing runs of text. Each
is fitted in \nm{clrs.g1.fit.lo} to \nm{clrs.g1.fit.hi} on one laptop core and runs cold in at most
\nm{clrs.g1.loops} loops at 64 nodes. The other fifteen read one quantity we supply, named in the table: the sum of a
distance and an edge weight for the shortest-path algorithms, the orientation of three points for the geometry tasks,
the cost of a split for the two dynamic programmes, prefix sums for the maximum subarray, closures of the graph with a
node or an edge removed for articulation points and bridges, and the open nodes of a search for the depth-first family.

The fitted rules are conjunctions of one to seven conditions, chosen by the fit from thousands of candidates, and they
read like the textbook (Box~1); for binary search and segment intersection, a learned score also decides among the
candidates that pass. With 10 training inputs instead of 100, \nm{few.exact} of the \nm{few.tasks} tasks we tried
stay exact, fitted in \nm{few.fit.lo} to \nm{few.fit.hi} milliseconds. Against the looped blocks of our earlier paper,
which trained for about an hour on a GPU per task\cite{mohsenvand2026zero}, the gain is in cost and in reach, and on
longest common subsequence, both string matchers, activity selection, task scheduling and Prim, in exactness.

\begin{figure}[tbp]\small
\fbox{\parbox{0.97\linewidth}{\textbf{Box 1 $|$ Six rules as fitted}, written out from the fitted tables; $k$ ranges over
third items, and $\notyet$ marks an item the partial answer has not reached.\\[0.4em]
\textbf{Insertion sort}, $u$ comes just before $v$: \ $\mathrm{key}_u<\mathrm{key}_v$, and no $k$ with
$\mathrm{key}_u<\mathrm{key}_k<\mathrm{key}_v$.\\[0.25em]
\textbf{Quickselect}, $i$ is the median: \ not $\#\{k:\mathrm{key}_k<\mathrm{key}_i\} < \#\{k:\mathrm{key}_k>\mathrm{key}_i\}$,
and not $\#\{k:\mathrm{key}_k\geq\mathrm{key}_i\} < \#\{k:\mathrm{key}_k<\mathrm{key}_i\}$.\\[0.25em]
\textbf{Kruskal}, $(u,v)$ is in the tree: \ an edge $u$--$v$ whose weight equals the bottleneck between $u$ and $v$, the
least possible largest weight on a path joining them (the cycle property).\\[0.25em]
\textbf{Activity selection}, $i$ is chosen: \ no chosen $k$ that finishes before $i$ finishes and after $i$ starts, the
answer read only downward along finish times.\\[0.25em]
\textbf{Breadth-first search}, $u$ is the parent of $v$: \ $\mathrm{depth}_u<\mathrm{depth}_v$, an edge $u\to v$, $u$
reached, $v$ not the source, and no $k$ with an edge $k\to v$, $\mathrm{depth}_k<\mathrm{depth}_v$ and a lower number
than $u$.\\[0.25em]
\textbf{Dijkstra}, with the sum supplied by hand (a ceiling): \ $\mathrm{dist}_u<\mathrm{dist}_v$, an edge $u\to v$, $u$
reached, $v$ not the source, and no $k$ with $\mathrm{dist}_k<\mathrm{dist}_v$ and
$\mathrm{dist}_k+w_{kv}<\mathrm{dist}_u+w_{uv}$.\\[0.4em]
Each rule is a conjunction of the pair's own comparisons and at most one witness test or two count comparisons, the
reads of the cold layer. The fit was told none of these conditions. The comparisons $\mathrm{depth}_u<\mathrm{depth}_v$
and $\mathrm{dist}_u<\mathrm{dist}_v$ are the level condition, which a finished answer cannot contradict. Read cold from $\notyet$, the
breadth-first rule places the nodes at distance $t$ from the source at loop $t$. All thirty rules are in Supplementary
Note~2.}}
\end{figure}

\begin{table}[tbp]\centering\scriptsize
\caption{\figlabel{Table 2}{All thirty CLRS-30 tasks fitted at the answer and run cold,} grouped by what we supply by
hand. Supplied: the quantity given by hand; results that rely on one are ceilings for the fitted check. Fit: wall clock on one laptop core. Loops: the most the cold run took on any
test input. Official 64: the library's score on the
official test. Exact to: the largest held-out size, in nodes or items, at which every held-out input was exact, with the
code frozen. Ours, looped: the better of the general and the cold looped blocks of our earlier paper; FloydNet: its
reported score\cite{yu2026floydnet}; both from Supplementary Table~4 of ref.~\cite{mohsenvand2026zero}. The averages mix
ceilings with the method and are not a like-for-like comparison.}\label{tab:clrs}
\setlength{\tabcolsep}{4.5pt}
\begin{tabular}{llrrrrrr}\toprule
Task & Supplied & Fit & Loops & Official 64 & Exact to & Ours, looped & FloydNet \\\midrule
\multicolumn{8}{l}{\textit{Comparisons and counts}} \\
Insertion sort & -- & 36 ms & 1 & 100.00 & 4,096 & 100.00 & 100.00 \\
Bubble sort & -- & 35 ms & 1 & 100.00 & 4,096 & 100.00 & 100.00 \\
Heapsort & -- & 34 ms & 1 & 100.00 & 4,096 & 100.00 & 100.00 \\
Quicksort & -- & 34 ms & 1 & 100.00 & 4,096 & 100.00 & 100.00 \\
Minimum & -- & 18 ms & 1 & 100.00 & 4,096 & 100.00 & 100.00 \\
Binary search & -- & 216 ms & 1 & 100.00 & 4,096 & 100.00 & 94.47 \\
Quickselect & -- & 18 ms & 1 & 100.00 & 4,096 & 100.00 & 80.33 \\
Activity selection & -- & 351 ms & 38 & 100.00 & 4,096 & 99.94 & 95.30 \\
Task scheduling & deadline $\div n$ & 479 ms & 8 & 100.00 & 4,096 & 99.86 & 91.05 \\
Breadth-first search & -- & 491 ms & 10 & 100.00 & 1,024 & 100.00 & 100.00 \\
Kruskal & -- & 94 ms & 1 & 100.00 & 1,024 & 100.00 & 99.63 \\
Prim & -- & 4.1 s & 87 & 100.00 & 1,024 & 99.98 & 99.67 \\
LCS length & index shift & 86 ms & 203 & 100.00 & 1,024 & 95.66 & 95.70 \\
Naive string matcher & -- & 212 ms & 3 & 100.00 & 5,120 & 99.07 & 100.00 \\
KMP matcher & -- & 207 ms & 3 & 100.00 & 5,120 & 99.15 & 99.43 \\
\addlinespace[2pt]
\multicolumn{8}{l}{\textit{One arithmetic primitive supplied}} \\
Bellman--Ford & sum & 3.9 s & 13 & 100.00 & 1,024 & 100.00 & 100.00 \\
Dijkstra & sum & 4.1 s & 13 & 100.00 & 1,024 & 100.00 & 97.67 \\
DAG shortest paths & sum & 4.2 s & 9 & 100.00 & 1,024 & 100.00 & 100.00 \\
Floyd--Warshall & sum & 43.3 s & 8 & 100.00 & 256 & 100.00 & 96.20 \\
Graham scan & orientation & 42 ms & 1 & 100.00 & 1,024 & 99.80 & 99.63 \\
Jarvis march & orientation & 42 ms & 1 & 100.00 & 1,024 & 99.92 & 99.50 \\
Segments intersect & orientation & 45 ms & 1 & 100.00 & fixed size & 99.33 & 99.90 \\
Matrix chain order & products, sums & 61 ms & 20 & 100.00 & 128 & 99.80 & 91.50 \\
Optimal BST & sums & 78 ms & 8 & 100.00 & 128 & 82.71 & 82.87 \\
\addlinespace[2pt]
\multicolumn{8}{l}{\textit{Algorithmic structure supplied}} \\
Maximum subarray & prefix sums & 458 ms & 1 & 100.00 & 4,096 & 87.86 & 86.20 \\
Articulation points & deletion closure & 23 ms & 1 & 100.00 & 1,024 & 100.00 & 100.00 \\
Bridges & deletion closure & 26 ms & 1 & 100.00 & 256 & 100.00 & 100.00 \\
Depth-first search & ``open'' & 47.5 s & 257 & 100.00 & 256 & 100.00 & 100.00 \\
Topological sort & ``open'', finish order & 47.6 s & 257 & 100.00 & 256 & 100.00 & 100.00 \\
Strong components & two passes & 47.5 s & 257 & 100.00 & 256 & 90.61 & 90.07 \\
\midrule
Average (30 tasks) & & & & 100.00 & & 98.46 & 96.64 \\
\bottomrule\end{tabular}

\end{table}

\textbf{One check for several algorithms.} The fitted check can be shared across tasks, as the looped block of our
earlier paper was. One check, told nothing about which task it serves, fitted jointly on breadth-first search, Prim,
Bellman--Ford, Dijkstra and shortest paths in an acyclic graph in \nm{shared.graph.fit}~s, is exact on all five, on the
official test and on fresh graphs of 1,024 nodes (Table~\ref{tab:shared}). One check serves the four sorts, fitted in
\nm{shared.sorts.fit}~s, and one the five selection tasks (activity selection, task scheduling, both convex hulls and
articulation points), with or without a bit that names the task. Each task still reads its own completion and supplied
quantity. Sharing has a limit: one check for all nine pointer tasks keeps the sorts and the shortest paths exact but
falls to \nm{shared.nine.bfs}\% on breadth-first search and \nm{shared.nine.prim}\% on Prim, with or without the task bit.

\begin{table}[tbp]\centering\scriptsize
\caption{\figlabel{Table 3}{One fitted check for several CLRS-30 tasks.} Each row is one check fitted jointly on the
training answers of all its tasks. Task named: whether a bit naming the task is among the inputs. Fit: wall clock on one
laptop core. Scores: the library's scorer on the official test and on fresh inputs of 1,024 nodes or items, per task.}
\label{tab:shared}
% Generated by gen/make_numbers.py from code/clrs/final/shared.json.
\begin{tabular}{lrlrp{0.2\linewidth}p{0.2\linewidth}}\toprule
One check for & Tasks & Task named & Fit & Official 64 (\%) & Fresh 1,024 (\%) \\\midrule
The four sorts & 4 & no & 0.6 s & all 100 & all 100 \\
Three shortest-path algorithms & 3 & no & 26.5 s & all 100 & all 100 \\
Five graph algorithms & 5 & no & 39.7 s & all 100 & all 100 \\
Five selection tasks & 5 & no & 4.2 s & all 100 & all 100 \\
Five selection tasks & 5 & yes & 3.8 s & all 100 & all 100 \\
Nine pointer tasks & 9 & no & 42.1 s & breadth-first search 31.6; Prim 35.9; others 100 & breadth-first search 24.9; Prim 3.1; others 100 \\
Nine pointer tasks & 9 & yes & 42.1 s & breadth-first search 31.6; Prim 35.9; others 100 & breadth-first search 24.9; Prim 3.1; others 100 \\
\bottomrule\end{tabular}

\end{table}

\subsection{Proofs at every size}\label{sec:proofs}

The fitted rules are small tables in which no size appears, so their correctness at every size reduces to finite
checks, which we ran by script (Supplementary Note~1 gives the arguments in full).

\emph{Sudoku.} A pair of cells has one of a fixed set of types, set by which coordinate relations agree, and the types
that a board realises do not change with the block size $b$ beyond $b=3$. A finite check confirms that the kept types
are exactly the row, column and box pairs at every $b$ from \nm{proof.b.lo} to \nm{proof.b.hi}, on all three fits. Every
kept cost is at least \nm{sudoku.wmin.lo}, no cost is negative, and the stopping threshold is \nm{sudoku.tau}, so a board
passes the check exactly when no row, column or box holds a value twice. A unit of $n$ cells with $n$ values then holds
each value once: the stop is sound at every size. Each removal of the cold stage follows from a count condition that
holds on every valid board, so the cold stage never removes a value of any solution, and with the complete search every
puzzle that has a solution is solved given enough moves.

\emph{Mazes.} On any tree maze the true path meets every kept range, so no removal is ever wrong. From the all-open
start, the cold stage removes cells with fewer than two possible path neighbours until what is left is a subtree whose
only leaves are the start and the goal, which is the path. No size and no edge term enter. A brute-force check over all
\nm{proof.mazes} spanning-tree mazes of the smallest lattices agrees. On Maze-Hard the clock is exactly the distance to
the goal, and the fitted rule takes the neighbour one step nearer and breaks ties by the label's own order, so it
returns the label on every maze of any size.

\emph{CLRS-30.} Proposition~\ref{prop:finite} applies to the sorts, minimum, Kruskal and breadth-first search, whose
rules use only own comparisons and witness tests. For the four sorts we checked every order of keys on up to
\nm{cert.sorts.items} items (\nm{cert.sorts.pairs} pairs each), and for minimum up to \nm{cert.min.items}. Kruskal's rule
uses only the pair's own conditions, and we checked every value its axioms allow against the cycle property. For
breadth-first search we took the depth as given, constrained only by the pairwise properties of a breadth-first depth,
and checked all \nm{cert.bfs.cases} pair cases on up to \nm{cert.bfs.items} items. These certify the check at the
answer at every size. That the cold run reaches the answer we argue, level by level, but have not machine-checked.

\subsection{What CRAM's released weights show}\label{sec:cramweights}

The same reading applies to CRAM's own released models, and it explains their two failures. Everything here comes from
the authors' public checkpoints and code, run unchanged.

\emph{The hidden layer cools.} For Maze-Hard, CRAM gives each open cell a hidden value $\psi$, updated as
$\psi_i \leftarrow \log(4+m) - \log \sum_{j \sim i} e^{-\psi_j}$ plus a source at the goal, which the authors' code
calls a soft-min Bellman backup. One learned number, the mass $m$, sets its temperature: a step of distance costs
$\log(4+m)$, so the temperature, in steps, is $1/\log(4+m)$. This is a walk distance, the family that runs from
random-walk distances at high temperature to the shortest-path distance at zero\cite{chebotarev2012walk,kivimaki2014developments}.
Over training the cost of a step rises from \nm{rel.step.init} to \nm{rel.step.lo} on seed 0, and to between
\nm{rel.step.lo} and \nm{rel.step.hi} on all three seeds, with the mass at \nm{rel.mass.end} and still rising at the
last step. Gradient descent is driving the layer towards a hard minimum, the read our cold layer takes, and a
temperature can only approach zero.

\emph{Where it is still warm, the tie count leaks in.} At any temperature above zero, the hidden value of a cell is
approximately its distance to the goal times the cost of a step, minus the logarithm of the number of shortest routes
from the cell. The walk-distance literature knows the second term as one that vanishes in the
limit\cite{francoisse2017bag}. Read cold, the distance and the tie count are separate numbers, the best witness and the
number tied with it in our earlier paper. Read warm, they are added, and a cell with many shortest routes looks nearer
than it is. The model's pointer adds learned scores for the four directions, \nm{rel.dir.down} for down, \nm{rel.dir.left}
for left, \nm{rel.dir.right} for right and \nm{rel.dir.up} for up, exactly the benchmark's tie-break order, so a large
enough room of equally short routes outweighs the gap between two directions. From the released numbers this predicts
the room size at which the model abandons the label's corridor for the room. We built \nm{rel.rooms.n} mazes of rooms
and corridors to test it and ran the full released models. On every seed,
\nm{rel.rooms.wrong} of the mazes are not given the label's path, from rooms of \nm{rel.rooms.first} by
\nm{rel.rooms.first} cells, as the law predicts, and \nm{rel.rooms.longer.lo} to \nm{rel.rooms.longer.hi} of those are given
a path two steps longer than the shortest (and one, on one seed, a set of cells that is not a path). Our rule, which
keeps the clock and the direction apart, is right on all of them.

\emph{The maze rule fails at the edge.} CRAM's Maze-Unique model sees a cell and its eight neighbours. Of the
\nm{rel.cert.interior} patterns a cell inside the board can see, the released weights follow the dead-end rule on all,
on every seed, which with a short induction proves the model exact inside the board at every size. On the board's edge
they break it on \nm{rel.cert.edge.bad} of \nm{rel.cert.edge} patterns, and in a corner on \nm{rel.cert.corner.bad} of
\nm{rel.cert.corner}. A wall next to a cell adds a cost for keeping it open, and beyond the edge there is no wall, so a
dead end that stops at the edge is never filled. CRAM's test set has no open cell on the edge. On our \nm{rel.pcd.n} tree
mazes with a ring of wall the released model solves \nm{rel.ring.solved} of \nm{rel.pcd.n} on every seed; on the same
mazes without the ring it solves none on any seed. The one Maze-Hard failure the authors report, a dead end of six
cells left open along the border of test maze 106, fits the same explanation.

\emph{CRAM's own stop can be proved sound.} The authors observe that their Sudoku search never stops at a wrong board.
Their models also hold \nm{rel.neg.lo} to \nm{rel.neg.hi} small negative costs on pairs of cells that share no unit, so
the observation does not follow at once. A short counting argument shows that a wrong board has at least two clashing
pairs. On each seed it therefore pays at least \nm{rel.stop.pay.lo} to \nm{rel.stop.pay.hi} in costs and can earn back
at most \nm{rel.stop.earn.lo} to \nm{rel.stop.earn.hi}, so every wrong board lies at least \nm{rel.stop.bound.lo}
above zero, far above the threshold of \nm{rel.stop.thr}. The stop is sound for every one of the $9^{81}$ boards.
